\documentclass[12pt,a4paper]{article}
\usepackage{amsmath, amssymb, amsthm}
\usepackage[utf8]{inputenc}
\usepackage[T1]{fontenc}
\usepackage{geometry}
\usepackage{enumitem}
\usepackage{xcolor}
\geometry{margin=2.5cm}

\newcommand{\easy}{{\color{green!60!black}$\bigstar$}}
\newcommand{\medium}{{\color{orange}$\bigstar\bigstar$}}
\newcommand{\hard}{{\color{red}$\bigstar\bigstar\bigstar$}}

\title{\textbf{Additional Exercises 3}\\
\large Power Series \& Radius of Convergence\\[4pt]
\normalsize Mirrors: SIT3 \quad|\quad
\easy\ Easy \quad \medium\ Medium \quad \hard\ Hard}
\date{}

\begin{document}
\maketitle
\thispagestyle{empty}

%----------------------------------------------------------
\section*{Exercise 1: Find the radius and domain of convergence}
%----------------------------------------------------------

\emph{Use the ratio test: $R = \lim |a_n/a_{n+1}|$, or root test: $R = 1/\limsup\sqrt[n]{|a_n|}$.
After finding $R$, check the endpoints $x = \pm R$ separately.}

\begin{enumerate}[leftmargin=*, label=(\alph*)]

\item \easy\quad $\displaystyle\sum_{n=1}^{\infty}\frac{x^n}{3^n \cdot n}$

\item \easy\quad $\displaystyle\sum_{n=0}^{\infty}\frac{x^n}{n!}$

\item \easy\quad $\displaystyle\sum_{n=1}^{\infty}(-1)^n\frac{x^n}{n}$

\item \easy\quad $\displaystyle\sum_{n=1}^{\infty}n!\cdot x^n$

\item \medium\quad $\displaystyle\sum_{n=1}^{\infty}\frac{(-1)^n\cdot x^{2n-1}}{5^n\sqrt{n+1}}$

\item \medium\quad $\displaystyle\sum_{n=1}^{\infty}\frac{2^n\cdot x^n}{(2n+1)^2\cdot\sqrt[4]{n}}$

\item \medium\quad $\displaystyle\sum_{n=1}^{\infty}\frac{(-1)^{n-1}\cdot x^{2n-2}}{3^{n-1}\cdot n\sqrt{n}}$

\item \medium\quad $\displaystyle\sum_{n=1}^{\infty}(-1)^n\ln(n)\cdot\frac{(2x+3)^n}{n\cdot(2n)!}$

\item \hard\quad $\displaystyle\sum_{n=1}^{\infty}\frac{1\cdot 3\cdot 5\cdots(2n-1)}{2\cdot 4\cdot 6\cdots(2n)}\cdot x^n$

\item \hard\quad $\displaystyle\sum_{n=1}^{\infty}\frac{(2n)!\cdot x^n}{2^{2n}(n!)^2}$

\item \hard\quad $\displaystyle\sum_{n=1}^{\infty}\frac{(3n-1)\cdot(-1)^n\cdot x^n}{n^3}$

\end{enumerate}

%----------------------------------------------------------
\section*{Exercise 2: Taylor series around $x=a$ --- find and identify domain}
%----------------------------------------------------------

\begin{enumerate}[leftmargin=*, label=(\alph*)]

\item \easy\quad $f(x)=e^x$ around $a=0$

\item \easy\quad $f(x)=\sin x$ around $a=0$

\item \easy\quad $f(x)=\dfrac{1}{1-x}$ around $a=0$

\item \medium\quad $f(x)=\dfrac{1}{x-5}$ around $a=1$

\item \medium\quad $f(x)=\sin x$ around $a=\pi/2$

\item \medium\quad $f(x)=\ln x$ around $a=1$

\item \hard\quad $f(x)=\dfrac{1}{(x-5)^3}$ around $a=1$
\quad\emph{(differentiate the series for $1/(x-5)$ term by term)}

\item \hard\quad $f(x)=\dfrac{1}{(x-2)^2}$ around $a=-3$

\end{enumerate}

%----------------------------------------------------------
\section*{Exercise 3: Expand into Maclaurin series}
%----------------------------------------------------------

\emph{Use known series: $e^u$, $\sin u$, $\cos u$, $\ln(1+u)$, $\frac{1}{1-u}$
with substitution.}

\begin{enumerate}[leftmargin=*, label=(\alph*)]

\item \easy\quad $f(x) = e^{-x^2}$

\item \easy\quad $f(x) = \sin(x^3)$

\item \medium\quad $f(x) = \dfrac{1}{x^2+4x+5}$
\quad\emph{(complete the square)}

\item \medium\quad $f(x) = xe^x$

\item \medium\quad $f(x) = \ln(1+x^2)$

\item \hard\quad $f(x) = \arctan\!\left(\dfrac{1+x^2}{1-x^2}\right)$

\item \hard\quad $f(x) = \dfrac{2x^4+3x^3-15x^2-5x+19}{x^2+x-6}$
\quad\emph{(polynomial division + partial fractions)}

\end{enumerate}

%----------------------------------------------------------
\section*{Exercise 4: High-order derivatives via Taylor series}
%----------------------------------------------------------

\begin{enumerate}[leftmargin=*, label=(\alph*)]

\item \medium\quad Compute $\dfrac{d^{17}}{dx^{17}}\!\left[\dfrac{1}{x-5}\right]\bigg|_{x=1}$

\item \medium\quad Compute $\dfrac{d^{10}}{dx^{10}}\!\left[\dfrac{1}{1-x}\right]\bigg|_{x=0}$

\item \hard\quad Given $f(x)$ analytic around $a$ with radius $R_a>1$.
Find $\displaystyle\sum_{n=0}^{\infty}\frac{f^{(n)}(a)}{n!}(1-a)^n$.

\end{enumerate}

\bigskip
\begin{center}
\small\itshape
\easy\ Ratio test, known Maclaurin series\quad
\medium\ Endpoint analysis, substitution into known series\quad
\hard\ Factorial products, term-by-term differentiation, abstract series.
\end{center}

\end{document}
